\chapter{Conclusion}
For the SABLE mid-stage scope, the implemented prototype addresses refactoring-induced misbinding of a registered S3/IAM security obligation. It checks the trusted baseline, normalizes a bounded Terraform/HCL subset, identifies a supported successor, evaluates the obligation on that successor, and publishes candidate-bound evidence to the shared gate. Controlled tests distinguish \texttt{PRESERVED}, \texttt{REGRESSED}, and \texttt{UNKNOWN}. The correction-containing local selection reported 27 tests passed and 10 deselected; the earlier CI run on commit \texttt{9a4dc3d} reported 44 focused tests passed. Current shared-workflow runs on branch head \texttt{753a55b} recorded \texttt{PRESERVED}/\texttt{REVIEW}, \texttt{REGRESSED}/\texttt{BLOCK}, and \texttt{UNKNOWN}/\texttt{REVIEW}. These are selected correctness checks rather than a full-suite or live-cloud result.

The principal design choice is to keep authorization validity separate from logical-asset correspondence. The focused audit also aligned the standalone evaluator's treatment of unresolved role bindings with the shared evaluator. That correction is commit \texttt{58b0a3d}, included in current branch head \texttt{753a55b}. The current boundaries remain the supported Terraform and IAM/S3 subset, synthetic authored scenarios, and the absence of a completed comparative research evaluation. Future SABLE work is listed in Chapter 10 as proposed research, not completed implementation.

This SABLE conclusion does not supply the missing team-wide research synthesis, confirmed panel action register, contribution records for other modules, administrative confirmations, or complete schedule and meeting evidence. The team must complete those report sections using authoritative records before submission.
